% Fichier généré automatiquement par tools/generate_corrections.py.
% Source : DYN/DYN-05-Methode/03_TT/03_TT.tex
% Statut : complet (2/2 question(s) renseignée(s), 0 reprise(s)).
\begin{corrigebox}[Corrigé extrait de la banque]
\CorrigeQuestion{1}
\begin{marginfigure}
\includegraphics[width=.6\linewidth]{03_TT_01_cor}
\end{marginfigure}
\CorrigeQuestion{2}
C'est une chaîne ouverte. On isole l'extrémité et on applique le théorème correspondant la mobilité : 
\begin{itemize}
\item on isole \textbf{2} et on réalise le théorème de la résultante dynamique en projection sur $\vj{0}$;
\item on isole \textbf{1+2} et on réalise le théorème de la résultante dynamique en projection sur $\vi{0}$.
\end{itemize}
\end{corrigebox}
